\documentclass[12pt]{article}
\usepackage{amsmath,amssymb,amsthm}
\usepackage[margin=1in]{geometry}

\newtheorem{theorem}{Theorem}
\newtheorem{lemma}[theorem]{Lemma}

\title{Problem 443: GCD Sequence}
\author{}
\date{}

\begin{document}
\maketitle

\section{Problem Statement}
Define the sequence $g(n)$ by $g(1) = 1$ and
\[
g(n) = g(n-1) + \gcd(n,\, g(n-1)) \quad \text{for } n \geq 2.
\]
Find $g(10^{15})$.

\section{Mathematical Foundation}

\begin{theorem}[Monotonicity]
The sequence $\{g(n)\}$ is strictly increasing with $g(n) \geq n$ for all $n \geq 1$.
\end{theorem}

\begin{proof}
By induction. Base case: $g(1) = 1 \geq 1$. Inductive step: $g(n) = g(n-1) + \gcd(n, g(n-1)) \geq g(n-1) + 1 > g(n-1)$, establishing strict monotonicity. Also $g(n) \geq (n-1) + 1 = n$.
\end{proof}

\begin{theorem}[Prime Generation --- Rowland, 2008]
If $g(n) - g(n-1) > 1$, then $g(n) - g(n-1)$ is prime for all $n \geq 5$.
\end{theorem}

\begin{proof}
Let $d = \gcd(n, g(n-1))$ with $d > 1$. Write $g(n-1) = n + \delta$ where $\delta = g(n-1) - n$. Then $d = \gcd(n, n + \delta) = \gcd(n, \delta)$. Since $\gcd(n, n-1) = 1$, the value $d$ divides $\delta$ but is coprime to consecutive integers near~$n$. A careful analysis of the deficit evolution shows that when $d > 1$, the value $d$ equals the smallest prime factor of $\gcd(n, \delta)$, which is itself prime. See Rowland (2008) and Cloitre (2011) for the complete argument.
\end{proof}

\begin{lemma}[Deficit Tracking]
\label{lem:deficit}
Define $\delta(n) = g(n) - n$. Then:
\begin{itemize}
\item If $\gcd(n, g(n-1)) = 1$: $\delta(n) = \delta(n-1)$.
\item If $\gcd(n, g(n-1)) = d > 1$: $\delta(n) = \delta(n-1) + d - 1$.
\end{itemize}
\end{lemma}

\begin{proof}
$\delta(n) = g(n) - n = g(n-1) + d - n = (g(n-1) - (n-1)) + (d - 1) = \delta(n-1) + (d - 1)$.
\end{proof}

\begin{theorem}[Skip Optimization]
When $\gcd(n, g(n-1)) = 1$ for consecutive values of~$n$, the sequence satisfies $g(n) = g(n-1) + 1$, allowing long stretches to be skipped. The next nontrivial GCD occurs at the smallest $n' > n$ such that $\gcd(n', \delta) > 1$, where $\delta = g(n-1) - (n-1)$ is the current deficit.
\end{theorem}

\begin{proof}
When $d = 1$, we have $g(n) = g(n-1) + 1$ and $\delta$ is constant. For future index $n' > n$, $\gcd(n', g(n'-1)) = \gcd(n', n' - 1 + \delta) = \gcd(n', \delta - 1 + n') = \gcd(n', \delta)$ (since $\gcd(n', n'-1) = 1$ and $g(n'-1) = n' - 1 + \delta$). Thus $d > 1$ iff $\gcd(n', \delta) > 1$, which occurs at the next multiple of any prime factor of~$\delta$.
\end{proof}

\section{Algorithm}

\begin{enumerate}
\item Initialize $g = 1$, $n = 2$.
\item At each step, compute $d = \gcd(n, g)$.
\item If $d = 1$: compute $\delta = g - n + 1$. Find the smallest prime factor~$p$ of~$\delta$, then skip to the next multiple of~$p$ at or above~$n$, advancing~$g$ linearly.
\item If $d > 1$: set $g \leftarrow g + d$, $n \leftarrow n + 1$.
\item Repeat until $n > N$.
\end{enumerate}

\section{Complexity Analysis}

\begin{itemize}
\item \textbf{Time:} Approximately $O(N / \log N)$ steps due to skip optimization. Each step requires a GCD computation ($O(\log N)$) and a smallest-prime-factor lookup.
\item \textbf{Space:} $O(\sqrt{N})$ for a prime sieve used in factoring the deficit, or $O(1)$ with trial division.
\end{itemize}

\section{Answer}

\[
\boxed{2744233049300770}
\]

\end{document}
